% ============================================================
% Chapter 3 — Methodology
% Thesis: Retrieval-Augmented Generation over Machine Learning
%         Knowledge Graphs
% Author: Yiğit Alp Öztorun
% Institution: KU Leuven
%
% Writing conventions:
%  - All verified numbers are marked with no annotation.
%  - All numbers awaiting final verification are marked with \todo{verify}.
%  - Figure and table \label commands use the plan's LaTeX labels.
%  - \todo{} is a placeholder — define it in your preamble, e.g.:
%      \newcommand{\todo}[1]{\textbf{[TODO: #1]}}
% ============================================================

\chapter{Methodology}
\label{ch:methodology}

% ============================================================
\section{Methodological Overview}
\label{sec:overview}
% ============================================================

This chapter presents the methodology of the proposed Retrieval-Augmented
Generation system designed for natural-language question answering over the
MLSea machine-learning knowledge graph.
The pipeline is organized into three successive and logically dependent stages.
In the \emph{pre-retrieval} stage, raw RDF entity data is extracted from a local
N-Triples file, assembled into structured canonical entity records, serialized
into textual representation-specific chunks, and embedded into a persistent
vector store.
In the \emph{retrieval} stage, each natural-language question is embedded with
the same model, and a suite of six candidate-ranking methods --- ranging from a
pure dense semantic baseline to hybrid methods that combine semantic similarity
with symbolic metadata signals --- is applied to produce a ranked top-10
candidate list per question.
In the \emph{post-retrieval} stage, the top candidates are filtered by retrieval
confidence, re-ranked using a cross-encoder, assembled into a structured evidence
context, and passed to a language model for answer generation.
The entire pipeline operates offline: all entity data is read directly from a
local copy of the MLSea knowledge graph, without dependency on a live SPARQL
endpoint or a running graph database.
This design ensures full reproducibility and experimental control.
The pipeline was evaluated on 265 answerable questions drawn from a manually
curated set of 280 questions, with questions targeting entities across three
entity types: scientific papers, benchmark datasets, and machine-learning models.
Figure~\ref{fig:pipeline_overview} illustrates the end-to-end flow of the
proposed system.

\begin{figure}[htbp]
  \centering
  % TODO: Insert pipeline overview figure (3-stage flow diagram)
  \fbox{\parbox{0.9\textwidth}{\centering\vspace{2cm}
    [Figure: End-to-end KG-RAG pipeline — three-stage flow diagram]
  \vspace{2cm}}}
  \caption{The end-to-end KG-RAG pipeline proposed in this thesis. Raw
    MLSea RDF data is parsed locally and entities are embedded into a
    vector store during the pre-retrieval stage. The retrieval stage ranks
    candidate entities per question using dense and hybrid methods. The
    post-retrieval stage re-ranks candidates, assembles a context block,
    and generates a natural-language answer.}
  \label{fig:pipeline_overview}
\end{figure}

% ============================================================
\section{Dataset and Knowledge Graph Source}
\label{sec:dataset}
% ============================================================

% -----------------------------------------------------------
\subsection{The MLSea Knowledge Graph}
\label{subsec:mlsea}
% -----------------------------------------------------------

The knowledge graph used in this thesis is MLSea~\todo{cite MLSea paper}, an
RDF representation of the machine-learning research landscape derived from the
Papers with Code~\todo{cite PwC} dataset.
MLSea encodes scientific papers, benchmark datasets, and machine-learning models
together with their semantic relationships.
The knowledge graph is available as a single N-Triples file,
\texttt{data/raw/pwc\_1.nt}, containing 26{,}606{,}202 RDF triples and
occupying 6.4~GB on disk.
In N-Triples format, each line encodes one atomic fact as a
\texttt{<subject> <predicate> <object> .} triple.

Three categories of entities are distinguished within the graph by their
IRI prefix:
\begin{itemize}
  \item \textbf{Scientific papers:} \texttt{http://w3id.org/mlsea/pwc/scientificWork/}
  \item \textbf{Benchmark datasets:} \texttt{http://w3id.org/mlsea/pwc/dataset/}
  \item \textbf{Machine-learning models:} \texttt{http://w3id.org/mlsea/pwc/model/}
\end{itemize}

% -----------------------------------------------------------
\subsection{Entity Types and RDF Structure}
\label{subsec:entity_types}
% -----------------------------------------------------------

The three entity types differ substantially in metadata richness.
Table~\ref{tab:entity_types} summarizes the entity counts and metadata
completeness observed after extraction.

\begin{table}[htbp]
  \centering
  \caption{MLSea entity types, their IRI prefixes, scale, and metadata
    completeness. The severe sparsity of dataset and model descriptions
    motivates distinct representation strategies for each entity type.}
  \label{tab:entity_types}
  \begin{tabular}{llrrrl}
    \toprule
    Entity type & IRI suffix & Total & Has title & Has description & Key predicates \\
    \midrule
    Paper   & \texttt{scientificWork/} & 489{,}584 & 92.1\% & 88.9\% (abstract)
            & \texttt{dcterms:title}, \texttt{fabio:abstract},
              \texttt{dcterms:creator}, \texttt{mlso:hasTaskType} \\
    Dataset & \texttt{dataset/}        & 12{,}651  & 100\%  & 0\%
            & \texttt{rdfs:label}, \texttt{dcat:keyword},
              \texttt{mlso:hasTaskType} \\
    Model   & \texttt{model/}          & 30{,}564  & 100\%  & 0\%
            & \texttt{rdfs:label}, \texttt{schema:codeRepository},
              \texttt{mls:realizes} \\
    \bottomrule
  \end{tabular}
\end{table}

Papers are the most metadata-rich entity type: they carry title and abstract
text, author lists, task associations, dataset links, method names, evaluation
metrics, and code repository links encoded via standard vocabularies including
\texttt{dcterms}, \texttt{fabio}, \texttt{dcat}, \texttt{mlso}, and
\texttt{mls}.
Datasets carry only a label and optional keyword and task associations; no
dataset in MLSea has a textual description.
Models similarly carry only a label and repository link; all semantic richness
for models comes from their linked-entity relationships within the graph.

% -----------------------------------------------------------
\subsection{Machine-Learning Domain Metadata in MLSea}
\label{subsec:mlsea_metadata}
% -----------------------------------------------------------

Raw RDF triples are not directly suitable as retrieval units.
A single scientific paper is described by tens or hundreds of triples, each
encoding one atomic fact with subject, predicate, and object IRIs.
Linked entities such as tasks, methods, datasets, and authors are referenced
by their IRIs, not by their human-readable labels.
To retrieve a label for a linked entity, the graph must be traversed: the
object IRI of one triple becomes the subject of another triple whose predicate
is \texttt{rdfs:label} or an equivalent.
Constructing a meaningful, self-contained text representation of an entity
therefore requires (i) collecting all triples with the entity as subject,
(ii) following linked IRIs to resolve their labels, and (iii) assembling the
resolved values into a coherent text.
This assembly process constitutes the core of the pre-retrieval stage and is
described in Section~\ref{sec:preretrieval}.

% ============================================================
\section{Question Set and Retrieval Task Definition}
\label{sec:questions}
% ============================================================

% -----------------------------------------------------------
\subsection{Question Dataset}
\label{subsec:question_dataset}
% -----------------------------------------------------------

The evaluation question set consists of 280 manually curated natural-language
questions targeting entities in the MLSea knowledge graph
(\texttt{data/questions/ml\_questions\_dataset.json}).
Each question is a realistic information-seeking query of the type a
machine-learning researcher might pose, such as asking for the paper that
introduced a specific model, the tasks associated with a benchmark dataset,
or the authors of a given publication.

Of the 280 questions, 265 are answerable: they have a valid gold entity IRI
(\texttt{target\_entity\_iri}) pointing to an entity present in the knowledge
graph.
The remaining 15 questions are explicitly marked as unanswerable
(\texttt{is\_answerable=false}) and are excluded from all quantitative
evaluation; they are retained in the dataset for completeness.

Each question record contains the fields listed in
Table~\ref{tab:question_dataset}.

\begin{table}[htbp]
  \centering
  \caption{Schema of the evaluation question dataset. All 280 questions
    target entities in the MLSea knowledge graph. Fifteen unanswerable
    questions are excluded from metric computation.}
  \label{tab:question_dataset}
  \begin{tabular}{lll}
    \toprule
    Field & Type & Description \\
    \midrule
    \texttt{id}               & string  & Unique question identifier \\
    \texttt{question}         & string  & Natural-language question text \\
    \texttt{question\_type}   & string  & Semantic type from the 35-type taxonomy \\
    \texttt{target\_entity\_iri} & string & Gold entity IRI in MLSea \\
    \texttt{answer}           & string/list & Structured gold answer \\
    \texttt{text\_answer}     & string  & Free-text gold answer \\
    \texttt{difficulty}       & string  & Difficulty label (Easy / Medium / Hard) \\
    \texttt{is\_answerable}   & boolean & Whether a gold entity exists \\
    \bottomrule
  \end{tabular}
\end{table}

The dataset covers 35 distinct question types, including paper-centric types
(e.g., \texttt{paper\_to\_tasks}, \texttt{paper\_to\_authors},
\texttt{paper\_description\_to\_title}), dataset-centric types
(e.g., \texttt{dataset\_to\_tasks}), model-centric types
(e.g., \texttt{repository\_to\_model}, \texttt{model\_family\_variant}),
and unanswerable variants.
Questions are distributed across three target entity types: 178 target
scientific papers, 57 target datasets, and 30 target models.

Difficulty labels are assigned as follows: Easy (40 questions), Medium (103
questions), Hard (\todo{verify exact count, either 77 or 87} questions),
and a subset with no difficulty label
(\todo{verify exact count, either 45 or 50} questions).

% -----------------------------------------------------------
\subsection{Gold Target Entity Design}
\label{subsec:gold_targets}
% -----------------------------------------------------------

Each answerable question's \texttt{target\_entity\_iri} uniquely identifies
the correct MLSea entity.
The entity type of the gold target is directly inferable from the IRI prefix:
IRIs containing \texttt{scientificWork/} identify papers, those containing
\texttt{dataset/} identify datasets, and those containing \texttt{model/}
identify models.
This allows the retrieval evaluation to automatically determine the expected
entity type for each question.

Across the 265 answerable questions, 37 unique gold entity IRIs are present.
Of these, 25 are paper entities and 12 are dataset entities; no questions
target model entities as gold, since model questions are answered by
identifying the correct paper or repository.
The 25 paper gold target IRIs are guaranteed to appear in the 200{,}000-paper
corpus subset used in this thesis (see Section~\ref{subsec:corpus_curation}).

% -----------------------------------------------------------
\subsection{Retrieval Task Formulation}
\label{subsec:retrieval_task}
% -----------------------------------------------------------

The retrieval task is defined as follows: given a natural-language question
$q$, the system must return a ranked list of at most $k = 10$ entity candidates
from the indexed corpus, such that the gold entity is ranked as highly as
possible.
Evaluation is performed using the following metrics, computed per question and
averaged across all 265 answerable questions:

\begin{itemize}
  \item \textbf{NDCG@10} (primary metric): Normalized Discounted Cumulative
        Gain at rank cutoff 10. If the gold entity is found at rank $r$,
        $\text{NDCG} = 1/\log_2(r + 1)$; otherwise 0.
  \item \textbf{Hit@}$k$ (for $k \in \{1, 5, 10\}$): Binary indicator
        of whether the gold entity appears within the top $k$ candidates.
  \item \textbf{MRR}: Mean Reciprocal Rank; $\text{MRR} = 1/r$ if the gold
        entity is found, otherwise 0.
\end{itemize}

NDCG@10 is selected as the primary metric because it rewards both retrieving
the correct entity and retrieving it at a high rank, which is critical for the
downstream post-retrieval stage.
The retrieval task is entity retrieval, not answer generation; answer-level
evaluation is performed separately in the post-retrieval stage.

The design of the corpus and the representation strategies described in
Section~\ref{sec:preretrieval} is directly constrained by the requirements of
this retrieval task.

% ============================================================
\section{Pre-Retrieval Phase: Entity Representation Construction}
\label{sec:preretrieval}
% ============================================================

% -----------------------------------------------------------
\subsection{Motivation for Pre-Retrieval Representations}
\label{subsec:preretrieval_motivation}
% -----------------------------------------------------------

Dense retrieval requires each entity to be represented as a fixed-size
embedding vector.
Raw RDF triples cannot be directly embedded because they are atomic,
disconnected, and encode IRIs rather than human-readable text.
A principled transformation step is required: entity-relevant triples must
be collected, linked IRI labels resolved, a readable text constructed, and
the text embedded.

Beyond this basic requirement, the thesis explores a more specific hypothesis:
different information types within an entity record --- title text, abstract,
structured metadata, graph-linked entities --- may carry different retrieval
signal depending on the question type and entity type.
This motivates constructing multiple representation strategies per entity type
and evaluating each independently, rather than defaulting to a single fixed
representation.
The pre-retrieval stage therefore serves two functions: preparing the corpus
for retrieval and empirically identifying the most effective representation
strategy per entity type.

% -----------------------------------------------------------
\subsection{RDF Extraction and Linked-Entity Label Resolution}
\label{subsec:rdf_extraction}
% -----------------------------------------------------------

Entity records are constructed by parsing the N-Triples source file using
a custom streaming parser.
Rather than loading the full graph into memory, the parser processes the
26.6M-triple file line by line using a regular-expression pattern that
extracts subject, predicate, and object from each N-Triples line.
This approach is memory-efficient and avoids dependency on a graph database
at extraction time.

The extraction proceeds in two passes over the source file.
In the first pass, all triples whose subject IRI matches the target entity
prefix are collected.
For paper extraction, this means retaining every triple whose subject begins
with \texttt{http://w3id.org/mlsea/pwc/scientificWork/}; analogous filtering
is applied for datasets and models.
Each unique subject IRI accumulates its associated predicate-object pairs into
a growing entity record.

The first pass produces records with linked entity objects represented as
raw IRIs rather than readable labels.
For example, a task associated with a paper is encoded as an IRI such as
\texttt{http://w3id.org/mlsea/pwc/task/ImageClassification}, not as the
string \textit{Image Classification}.
A second pass over the source file resolves these linked IRIs to human-readable
labels.
For each encountered object IRI, the parser searches for a label predicate in
the following order of preference: \texttt{dcterms:title}, \texttt{rdfs:label},
\texttt{foaf:name}, \texttt{schema:name}.
If none of these predicates is found for an IRI, the local name is extracted
from the IRI (the fragment after \texttt{\#} or the final path segment) and
used as a fallback label.

Figure~\ref{fig:rdf_to_chunk} illustrates the transformation from raw RDF
triples through the two-pass extraction to the final text chunk.

\begin{figure}[htbp]
  \centering
  % TODO: Insert RDF triples → canonical record → chunk transformation diagram
  \fbox{\parbox{0.9\textwidth}{\centering\vspace{2cm}
    [Figure: Two-pass RDF extraction and label resolution]
  \vspace{2cm}}}
  \caption{Transformation of raw RDF triples into an entity-centric text chunk.
    Pass~1 collects all triples for a target entity by subject IRI prefix
    matching. Pass~2 re-scans the file to resolve linked IRI labels. The
    assembled canonical entity record is then serialized into a
    representation-specific chunk whose \texttt{source\_text} field is
    embedded for retrieval.}
  \label{fig:rdf_to_chunk}
\end{figure}

% -----------------------------------------------------------
\subsection{Canonical Entity Record Construction}
\label{subsec:canonical_records}
% -----------------------------------------------------------

The output of the two-pass extraction is a \emph{canonical entity record} for
each entity --- a structured dictionary that consolidates all resolved metadata
into a single in-memory object.
This record is the intermediate unit between raw triples and the
representation-specific text chunks produced later; the two must not be
conflated.

For scientific papers, the canonical record includes the following fields:
\texttt{paper\_id}, \texttt{title}, \texttt{abstract}, \texttt{authors},
\texttt{keywords}, \texttt{tasks}, \texttt{datasets}, \texttt{methods},
\texttt{metrics}, \texttt{implementations}, \texttt{publication\_year},
and \texttt{linked\_entities} (a list of objects carrying predicate label,
object label, object IRI, and inferred entity type).
For datasets, the record includes the dataset IRI, label, keywords, tasks,
and linked paper IRIs.
For models, the record includes the model IRI, label, repository names and
URLs, and linked entities grouped by predicate.

Canonical records are serialized to JSONL files:
\texttt{papers\_master.jsonl} (489{,}584 records),
\texttt{datasets\_master.jsonl} (12{,}651 records), and
\texttt{models\_master.jsonl} (30{,}564 records).
These files serve as the source of truth for all downstream stages, including
post-retrieval context assembly.

% -----------------------------------------------------------
\subsection{Corpus Curation and Subset Selection}
\label{subsec:corpus_curation}
% -----------------------------------------------------------

The full paper corpus of 489{,}584 records is computationally expensive to
embed and index within the scope of a thesis experiment.
A curated subset of 200{,}000 papers was therefore selected, reducing the
embedding and storage cost while maintaining evaluation validity.

The subset is constructed using a gold-first inclusion strategy: all 25 paper
gold targets identified from the question dataset are included unconditionally
and placed at the front of the subset.
The remaining 199{,}975 positions are filled sequentially from the master file.
This guarantees that no gold paper target is absent from the retrieval corpus,
which would artificially lower the evaluation ceiling.

Dataset and model entities are not subsetted.
All 12{,}651 datasets and 30{,}564 models are indexed in their entirety.
Of the 37 unique gold entity IRIs in the question dataset, 25 are paper
entities (all included in the subset) and 12 are dataset entities (all present
in the full dataset corpus).
The detailed corpus statistics are summarized in
Table~\ref{tab:corpus_stats}.

\begin{table}[htbp]
  \centering
  \caption{Corpus construction statistics for each entity type. Paper records
    are subsetted to 200{,}000 for computational feasibility; all 25 paper
    gold targets are guaranteed to be included. Dataset and model corpora are
    used in full.}
  \label{tab:corpus_stats}
  \begin{tabular}{lrrrp{5cm}}
    \toprule
    Entity type & Total extracted & Corpus used & Gold targets & Notes \\
    \midrule
    Paper   & 489{,}584 & 200{,}000 & 25 (all included) &
              Gold-first inclusion; remaining slots filled sequentially \\
    Dataset & 12{,}651  & 12{,}651  & 12 (all present)  &
              No subset applied \\
    Model   & 30{,}564  & 30{,}564  & 0                 &
              No model questions have model-entity gold targets \\
    \bottomrule
  \end{tabular}
\end{table}
  
% - - - - - - - - - - - - - - - - - - - - - - - - - - - - -
\subsubsection{Paper Representations}
\label{subsubsec:paper_representations}
% - - - - - - - - - - - - - - - - - - - - - - - - - - - - -

Six representation strategies are defined for scientific papers, spanning a
spectrum from minimal text content to graph-informed richness.

The \texttt{title\_only} and \texttt{abstract\_only} strategies serve as
ablation baselines that isolate the retrieval signal of a single field.
They establish whether the title or abstract alone is sufficient for the
retrieval task.
The \texttt{title\_abstract} strategy combines both primary text fields and
constitutes the natural text-based baseline.

The \texttt{predicate\_filtered} strategy extends the title and abstract with
structured metadata fields: up to five items each from tasks, datasets,
methods, metrics, and implementations.
These fields are populated by traversing named ML-relevant predicates,
ensuring that only semantically meaningful attributes are included.

The \texttt{enriched\_metadata} strategy builds on predicate-filtered content
by additionally incorporating authors (up to six) and implementations (up to
three), targeting question types that require knowledge of authorship or code
availability.

The \texttt{one\_hop} strategy uses the full one-hop neighborhood of the paper
in the knowledge graph, incorporating all linked entities grouped by predicate
category, up to a limit of twelve linked entities.
This representation captures graph-structural context that is not present in
the paper's direct textual fields.

% - - - - - - - - - - - - - - - - - - - - - - - - - - - - -
\subsubsection{Dataset Representations}
\label{subsubsec:dataset_representations}
% - - - - - - - - - - - - - - - - - - - - - - - - - - - - -

Four representation strategies are defined for benchmark datasets.
Datasets in MLSea are metadata-sparse: none has a textual description, and
only 54.8\% have keyword associations.
This sparsity motivates testing whether even minimal representations suffice
for retrieval.

The \texttt{dataset\_title\_only} strategy encodes only the dataset label.
The \texttt{dataset\_metadata} strategy appends publication year, keywords,
and task associations.
The \texttt{dataset\_enriched\_metadata} strategy additionally incorporates
related paper IRIs and implementation links.

The \texttt{dataset\_predicate\_filtered} strategy is distinctive in that it
applies a corpus-level filter: datasets with no keywords, task associations,
related papers, or related implementations are excluded from this collection
entirely, reducing the indexed corpus from 12{,}651 to 6{,}928 entities.
This approach tests whether a higher-quality, lower-coverage corpus produces
better retrieval precision despite the reduced recall ceiling.
 
% -----------------------------------------------------------
\subsection{Embedding Generation and Vector Indexing}
\label{subsec:embedding}
% -----------------------------------------------------------

Each chunk's \texttt{source\_text} is encoded as a dense vector using the
\texttt{sentence-transformers/all-MiniLM-L6-v2}
model~\todo{cite all-MiniLM-L6-v2}.
This model produces 384-dimensional embeddings with L2 normalization, and is
well-suited to sentence- and paragraph-length inputs.
The configuration is summarized in Table~\ref{tab:embedding_config}.

\begin{table}[htbp]
  \centering
  \caption{Configuration parameters for embedding generation and vector store
    indexing. All 14 representation corpora were embedded using the same model
    and indexed in separate ChromaDB collections.}
  \label{tab:embedding_config}
  \begin{tabular}{ll}
    \toprule
    Parameter & Value \\
    \midrule
    Embedding model    & \texttt{sentence-transformers/all-MiniLM-L6-v2} \\
    Embedding dimension & 384 \\
    Normalization      & L2-normalized \\
    Vector store       & ChromaDB (local persistent) \\
    Index type         & HNSW~\todo{cite Malkov \& Yashunin 2018} \\
    Distance metric    & Cosine similarity \\
    Collections        & 18 (one per representation strategy) \\
    Total storage      & 8.2~GB \\
    Batch size         & \todo{verify from config/pre\_retrieval\_config.json} \\
    \bottomrule
  \end{tabular}
\end{table}

Embeddings are generated in batches and stored in a persistent local
ChromaDB instance.
One collection is created per representation strategy, resulting in 18
collections in total: six for papers, four for datasets, four for models,
and four auxiliary collections used for evaluation bookkeeping.
ChromaDB uses an HNSW index internally for approximate nearest-neighbour
search.

During pre-retrieval evaluation, each question text is embedded using the
same model and the top-10 most similar entities are retrieved by cosine
similarity from the relevant collection.

% -----------------------------------------------------------
\subsection{Pre-Retrieval Output}
\label{subsec:preretrieval_output}
% -----------------------------------------------------------

For each of the 14 representation strategies, the pre-retrieval stage
produces a ranked top-10 candidate list per question, serialized to JSON.
The output path follows the pattern:
\begin{center}
\texttt{data/results/pre\_retrieval\_results/\{entity\_type\}\_results/\{representation\}/top10.json}
\end{center}

Each entry contains the question identifier, the gold entity IRI, and a
ranked list of up to ten documents, each carrying entity IRI, title, cosine
similarity score, and rank.

The pre-retrieval evaluation identifies the representation strategy with the
highest NDCG@10 per entity type, which is then designated as the primary
input to the retrieval stage.
The selected representations are:
\begin{itemize}
  \item \textbf{Papers:} \texttt{enriched\_metadata}
  \item \textbf{Datasets:} \texttt{dataset\_title\_only}
  \item \textbf{Models:} \texttt{model\_predicate\_filtered}
\end{itemize}

The retrieval stage operates exclusively on these pre-computed top-10 lists;
it does not re-embed questions or re-query the vector store.
Having established the best retrieval representations and produced top-10
candidate lists for each entity type, the retrieval stage applies a suite of
re-ranking and fusion strategies to refine these candidate sets.

% ============================================================
\section{Retrieval Phase: Candidate Entity Ranking}
\label{sec:retrieval}
% ============================================================

% -----------------------------------------------------------
\subsection{Retrieval Objective and Design Rationale}
\label{subsec:retrieval_objective}
% -----------------------------------------------------------

The retrieval stage takes the pre-computed top-10 candidate lists as input
and applies re-ranking, symbolic filtering, or cross-representation fusion
to produce a refined top-$k$ candidate list per question.
Its purpose is to test whether signals derived from the knowledge graph's
structure and from the question's metadata type can improve over the pure
semantic similarity signal already captured by the pre-retrieval stage.

Six retrieval methods are defined and evaluated.
Figure~\ref{fig:retrieval_methods} illustrates their composition, and
Table~\ref{tab:retrieval_methods} summarizes the signal types each method
employs.

\begin{figure}[htbp]
  \centering
  % TODO: Insert retrieval method taxonomy figure
  \fbox{\parbox{0.9\textwidth}{\centering\vspace{2cm}
    [Figure: Retrieval method hierarchy and composition]
  \vspace{2cm}}}
  \caption{The six retrieval methods evaluated in this thesis. The dense
    baseline uses pre-computed semantic scores directly. Hybrid methods add
    symbolic re-ranking layers. RRF methods fuse multiple representation lists
    before applying symbolic filters.}
  \label{fig:retrieval_methods}
\end{figure}

\begin{table}[htbp]
  \centering
  \caption{Design composition of the six retrieval methods. Each method
    applies a distinct combination of semantic similarity, symbolic filtering,
    and cross-representation fusion.}
  \label{tab:retrieval_methods}
  \begin{tabular}{lccccc}
    \toprule
    Method & Semantic & Type & Richness & Predicate & Fusion \\
           & base     & filter & boost  & boost     & (RRF)  \\
    \midrule
    \texttt{pure\_semantic\_dense}          & \checkmark & & & & \\
    \texttt{hybrid\_type\_filtering}        & \checkmark & \checkmark & & & \\
    \texttt{hybrid\_type\_onehop\_filtering}& \checkmark & \checkmark & \checkmark & & \\
    \texttt{hybrid\_predicate\_aware}       & \checkmark & & & \checkmark & \\
    \texttt{optional\_rrf\_fusion}          & \checkmark & & & & \checkmark \\
    \texttt{optional\_rrf\_symbolic}        & \checkmark & \checkmark & & \checkmark & \checkmark \\
    \bottomrule
  \end{tabular}
\end{table}

All methods output a ranked list of at most $k = 10$ candidates per question.

% -----------------------------------------------------------
\subsection{Dense Retrieval Baseline}
\label{subsec:dense_baseline}
% -----------------------------------------------------------

The \texttt{pure\_semantic\_dense} method serves as the retrieval baseline.
It passes the pre-retrieval top-10 lists through without modification: for
each question, the best-representation top-10 list for the expected entity
type is returned as-is, preserving the original cosine similarity scores and
ranking.
The entity type is determined from the \texttt{target\_entity\_iri} prefix.
This method establishes the retrieval performance achievable by bi-encoder
semantic similarity alone and provides a reference against which all
subsequent methods are compared.

% -----------------------------------------------------------
\subsection{Symbolic and Metadata-Aware Hybrid Retrieval}
\label{subsec:hybrid_retrieval}
% -----------------------------------------------------------

% - - - - - - - - - - - - - - - - - - - - - - - - - - - - -
\subsubsection{Type-Filtered Retrieval}
\label{subsubsec:type_filtering}
% - - - - - - - - - - - - - - - - - - - - - - - - - - - - -

The \texttt{hybrid\_type\_filtering} method adds a type-consistency
constraint to the dense baseline.
Each candidate's entity type is inferred from its IRI prefix and compared to
the expected entity type for the question.
Candidates of the wrong type are removed.
If all candidates are filtered out, the original dense list is restored
unchanged as a fallback to preserve coverage.

This filter is motivated by the observation that the bi-encoder retrieval
operates over a mixed corpus spanning all entity types.
Without type filtering, questions targeting papers may retrieve datasets or
models that happen to be semantically similar in the embedding space but are
not valid answers.

% - - - - - - - - - - - - - - - - - - - - - - - - - - - - -
\subsubsection{One-Hop Richness-Boosted Retrieval}
\label{subsubsec:onehop_filtering}
% - - - - - - - - - - - - - - - - - - - - - - - - - - - - -

The \texttt{hybrid\_type\_onehop\_filtering} method applies type filtering
and then re-ranks the surviving candidates by a soft graph-richness signal.
A richness score is computed for each candidate by counting the number of
non-empty metadata fields among \texttt{tasks}, \texttt{datasets},
\texttt{methods}, \texttt{metrics}, and \texttt{implementations}, with a
bonus of two points for the presence of linked-entity text in the
representation.

The richness signal is combined with the semantic score through a soft blend:
\[
  \text{score}_{\text{blended}} = \text{score}_{\text{semantic}}
    + \varepsilon \cdot \text{score}_{\text{richness}},
  \quad \varepsilon = 0.01.
\]

The small value of $\varepsilon$ ensures that the richness contribution is at
most $0.07$ (seven fields at maximum richness), so richness can only resolve
near-ties in semantic score and cannot override a clearly superior semantic
match.
This soft-blend design prioritizes retrieval precision while providing a
tie-breaking signal favoring entities with richer knowledge-graph metadata.

% - - - - - - - - - - - - - - - - - - - - - - - - - - - - -
\subsubsection{Predicate-Aware Retrieval}
\label{subsubsec:predicate_aware}
% - - - - - - - - - - - - - - - - - - - - - - - - - - - - -

The \texttt{hybrid\_predicate\_aware\_filtering} method applies
question-type-specific metadata boosting.
Rather than boosting all metadata-rich entities uniformly, it identifies
which metadata field is most likely to be discriminative for the question's
semantic type.

The boosting rules are derived from the question-type taxonomy:
\begin{itemize}
  \item Task-related question types (e.g., \texttt{paper\_to\_tasks},
        \texttt{paper\_by\_task\_pair}, \texttt{dataset\_to\_tasks}):
        candidates with a non-empty \texttt{tasks} field are boosted.
  \item Implementation question types (\texttt{paper\_to\_implementation}):
        candidates with a non-empty \texttt{implementations} field are boosted.
  \item Year question types (\texttt{paper\_to\_publication\_year},
        \texttt{dataset\_to\_publication\_year}): candidates with a known
        publication year are boosted.
  \item Repository and model family types (\texttt{repository\_to\_model},
        \texttt{model\_family\_variant}): candidates with linked-entity text
        are boosted.
  \item Keyword types (\texttt{paper\_to\_keywords}): candidates with a
        non-empty \texttt{keywords} field are boosted.
\end{itemize}

Boosted candidates are placed ahead of non-boosted candidates; the relative
ordering within each group is preserved.
This method aligns the candidate ranking with the specific informational
requirement expressed by the question's type, without modifying semantic
scores.

% -----------------------------------------------------------
\subsection{Multi-Representation Fusion}
\label{subsec:rrf}
% -----------------------------------------------------------

% - - - - - - - - - - - - - - - - - - - - - - - - - - - - -
\subsubsection{Reciprocal Rank Fusion}
\label{subsubsec:rrf_fusion}
% - - - - - - - - - - - - - - - - - - - - - - - - - - - - -

The \texttt{optional\_rrf\_fusion} method combines top-10 candidate lists
from multiple representation strategies for the same entity type using
Reciprocal Rank Fusion (RRF)~\todo{cite Cormack et al. 2009}.
The RRF score for a candidate at rank $r$ in a given list is:
\[
  \text{RRF}(r) = \frac{1}{k + r}, \quad k = 60.
\]
A candidate's total score is the sum of its RRF contributions across all
lists in which it appears.
Candidates that appear in multiple representation lists accumulate higher
total scores, reflecting consensus across different information views of the
same entity.

The fusion groups are:
\begin{itemize}
  \item \textbf{Papers:} \texttt{enriched\_metadata},
        \texttt{predicate\_filtered}, \texttt{one\_hop} (3 representations).
  \item \textbf{Datasets:} \texttt{dataset\_title\_only},
        \texttt{dataset\_enriched\_metadata} (2 representations).
  \item \textbf{Models:} \texttt{model\_predicate\_filtered},
        \texttt{model\_enriched\_metadata} (2 representations).
\end{itemize}

The RRF constant $k = 60$ is the standard value used in the literature and
prevents the fusion score from being dominated by the top-ranked position.

% - - - - - - - - - - - - - - - - - - - - - - - - - - - - -
\subsubsection{RRF with Symbolic Filtering}
\label{subsubsec:rrf_symbolic}
% - - - - - - - - - - - - - - - - - - - - - - - - - - - - -

The \texttt{optional\_rrf\_symbolic} method applies the full symbolic
signal stack on top of RRF fusion.
It proceeds as a three-stage cascade: (1) RRF fusion is performed across
the designated representation groups; (2) type filtering is applied to
the fused candidates; (3) predicate-aware boosting is applied to the
type-filtered candidates.
This combination tests whether the diversity benefit of fusion and the
precision benefit of symbolic signals compound positively when applied
jointly.

% -----------------------------------------------------------
\subsection{Top-$k$ Candidate Generation}
\label{subsec:topk}
% -----------------------------------------------------------

All six methods truncate their output to at most $k = 10$ candidates per
question.
Each candidate in the output carries the fields listed in
Table~\ref{tab:candidate_schema}.

\begin{table}[htbp]
  \centering
  \caption{Schema of a retrieval candidate record. Each question's result
    contains up to ten candidate records; the post-retrieval stage reads all
    listed fields to assemble evidence context.}
  \label{tab:candidate_schema}
  \begin{tabular}{lll}
    \toprule
    Field & Type & Description \\
    \midrule
    \texttt{entity\_id}           & string       & Original entity IRI \\
    \texttt{entity\_type}         & string       & \texttt{paper}, \texttt{dataset}, or \texttt{model} \\
    \texttt{title\_or\_label}     & string       & Entity title or label \\
    \texttt{original\_rank}       & integer      & Rank in the pre-retrieval top-10 \\
    \texttt{new\_rank}            & integer      & Rank after retrieval-stage re-ranking \\
    \texttt{original\_score}      & float        & Pre-retrieval cosine similarity score \\
    \texttt{method\_score}        & float        & Score assigned by the retrieval method \\
    \texttt{representation\_source} & string     & Representation strategy used \\
    \texttt{metadata.tasks}       & list[string] & Associated ML tasks \\
    \texttt{metadata.authors}     & list[string] & Author names (papers only) \\
    \texttt{metadata.keywords}    & list[string] & Keywords \\
    \texttt{metadata.implementations} & list[string] & Code repository links \\
    \texttt{metadata.source\_text}& string       & Full representation text \\
    \bottomrule
  \end{tabular}
\end{table}

For RRF methods, three additional fields are present:
\texttt{rrf\_score} (total RRF score),
\texttt{contributing\_representations} (list of representation strategies
that ranked this candidate), and
\texttt{original\_ranks\_by\_representation} (a mapping from representation
name to the rank assigned in that list).

% -----------------------------------------------------------
\subsection{Retrieval Output}
\label{subsec:retrieval_output}
% -----------------------------------------------------------

Results for all six methods are stored under
\texttt{data/results/retrieval/\{method\_name\}/}.
Per-question results include the question metadata, the top-10 candidate list,
and the per-question metric values (Hit@1, Hit@5, Hit@10, MRR, NDCG).
Aggregate summaries are segmented by entity type, difficulty, and question
type.

The top-$k$ candidate list, together with the canonical entity records in
\texttt{papers\_master.jsonl}, constitutes the input to the post-retrieval
stage.
The retrieval stage produces a ranked list of candidate entities per question,
each accompanied by its structured metadata and representation text.
The post-retrieval stage takes this candidate list as its evidence base,
assembles a context block, re-ranks candidates using a cross-encoder, and
generates a final natural-language answer.

% ============================================================
\section{Post-Retrieval Phase: Context Assembly and Answer Generation}
\label{sec:postretrieval}
% ============================================================

% -----------------------------------------------------------
\subsection{Post-Retrieval Objective and Design Rationale}
\label{subsec:postretrieval_objective}
% -----------------------------------------------------------

The retrieval stage provides a ranked list of up to ten candidate entities per
question, but this list is not a natural-language answer.
The post-retrieval stage has three responsibilities: (i) filter out
low-confidence candidates based on retrieval score, (ii) re-rank remaining
candidates using a cross-encoder that evaluates question-candidate relevance
jointly, and (iii) assemble the top candidates into a structured context block
for language-model-based answer generation.

The post-retrieval stage operates fully offline, reading canonical entity
records and pre-retrieval representation texts from local JSONL files without
querying the vector store.
Figure~\ref{fig:postretrieval_pipeline} provides an overview of the
post-retrieval pipeline.

\begin{figure}[htbp]
  \centering
  % TODO: Insert post-retrieval flow diagram
  \fbox{\parbox{0.9\textwidth}{\centering\vspace{2cm}
    [Figure: Post-retrieval context assembly and answer generation pipeline]
  \vspace{2cm}}}
  \caption{The post-retrieval pipeline. Retrieved candidates are first
    filtered by a hard retrieval-score threshold, then re-ranked by a
    cross-encoder. The top-3 candidates are assembled into a structured
    context block, which is provided to a language model to generate a
    natural-language answer.}
  \label{fig:postretrieval_pipeline}
\end{figure}

% -----------------------------------------------------------
\subsection{Candidate Context Construction}
\label{subsec:context_construction}
% -----------------------------------------------------------

For each candidate in the top-$k$ list, the post-retrieval stage resolves
the corresponding canonical entity record from the master JSONL file.
A structured context string is assembled for each candidate, containing the
following fields: \texttt{Title}, \texttt{Publication Year}, \texttt{Authors},
\texttt{Tasks}, \texttt{Keywords}, \texttt{Datasets}, \texttt{Methods},
\texttt{Metrics}, \texttt{Implementations}, \texttt{Abstract} (truncated to
1{,}200 characters), and \texttt{Retrieved Representation} (the pre-retrieval
\texttt{source\_text}, truncated to 1{,}400 characters).

This per-candidate context string is the primary unit used in cross-encoder
re-ranking and, for the selected top candidates, in language model
conditioning.

% -----------------------------------------------------------
\subsection{Evidence Selection and Re-Ranking}
\label{subsec:reranking}
% -----------------------------------------------------------

Evidence selection proceeds in two stages.

In the first stage, a hard confidence filter removes candidates whose
retrieval cosine score falls below 0.20.
If all candidates fall below this threshold, the question is flagged as
unanswerable from the available context, and a fixed unanswerable response
is returned without invoking subsequent stages.

In the second stage, the surviving candidates are re-ranked by a cross-encoder
model (\texttt{cross-encoder/ms-marco-MiniLM-L-6-v2}~\todo{cite}).
The cross-encoder takes the question and each candidate's context string as a
joint input pair and produces a relevance score.
Unlike the bi-encoder used in pre-retrieval, the cross-encoder attends to both
inputs simultaneously, enabling more accurate relevance estimation at the cost
of higher computational complexity.
The candidates are re-ordered by cross-encoder score, and the top-3 candidates
are selected as the final evidence set.

The rationale for two-stage filtering is efficiency: the hard score filter
reduces the candidate set inexpensively before invoking the cross-encoder,
which is applied only to a small, filtered candidate set.

% -----------------------------------------------------------
\subsection{Answer Generation}
\label{subsec:generation}
% -----------------------------------------------------------

The top-3 re-ranked candidates are serialized into a structured
\texttt{<CONTEXT>} block that lists, for each candidate, the entity IRI,
retrieval score, cross-encoder score, and the full context string.
This context block is provided to a language model as part of a structured
prompt.

The system prompt instructs the model to act as an offline MLSea RAG
assistant, to answer the question using only the provided context, and to
respond with a fixed unanswerable statement if the context is insufficient.
The generation model is \texttt{TinyLlama/TinyLlama-1.1B-Chat-v1.0}
~\todo{cite TinyLlama}, a small chat-tuned language model deployed locally.
Decoding uses a temperature of 0.1 for near-deterministic output, with a
maximum of 256 new tokens.

If the context block is empty after the hard filtering stage (no candidates
survived), the generation step is skipped and the unanswerable response is
returned directly.

% -----------------------------------------------------------
\subsection{Post-Retrieval Output and Evaluation}
\label{subsec:postretrieval_output}
% -----------------------------------------------------------

The post-retrieval stage produces, per question: the assembled context block,
the generated natural-language answer, and a suite of answer-quality metrics.
Three metrics are used to assess answer quality:

\begin{itemize}
  \item \textbf{SAS (Semantic Answer Similarity)~\todo{cite}:} the cosine
        similarity between the embedding of the generated answer and the
        embedding of the gold text answer.
        SAS captures semantic equivalence that surface-form metrics may miss.
  \item \textbf{ROUGE-L~\todo{cite Lin 2004}:} the longest common subsequence
        overlap between the generated answer and the gold answer.
        ROUGE-L measures lexical overlap as a complementary surface-level metric.
  \item \textbf{LLM-as-a-judge:} a second model pass that compares the
        generated answer to the ground truth and returns a binary correctness
        score (correct / incorrect) together with a reasoning trace.
        The judge responses are retained for manual audit.
\end{itemize}

Quantitative results and their interpretation are presented in Chapter~4.

% -----------------------------------------------------------
\section*{Chapter Summary}
% -----------------------------------------------------------

This chapter has described the full methodology of the proposed KG-RAG
pipeline over the MLSea knowledge graph.
Beginning with local extraction and linked-label resolution from
26{,}606{,}202 raw RDF triples, the pipeline constructs canonical entity
records for 489{,}584 scientific papers, 12{,}651 datasets, and 30{,}564
models.
Fourteen representation strategies --- six for papers, four for datasets, and
four for models --- are embedded into a 384-dimensional vector space using the
\texttt{sentence-transformers/all-MiniLM-L6-v2} model and indexed in a
persistent ChromaDB vector store across 18 collections.
Six retrieval methods, from a pure dense semantic baseline to Reciprocal Rank
Fusion combined with symbolic filtering, are evaluated on 265 answerable
questions targeting papers, datasets, and models.
The top candidates from retrieval are passed to a post-retrieval stage that
applies hard confidence filtering, cross-encoder re-ranking, and language
model generation.
Chapter~4 presents the quantitative evaluation results for all 14
representation strategies and six retrieval methods, including performance
analysis segmented by entity type, question difficulty, and question type.
